\scanpage{211}
% L1--4 (haut de la page)
that the automorphisms of such a $G_\nu$ ($\simeq A_{g,\nu}$) which fix the class $U_\nu$ i.e. which \unsure{transform|transport} $H_\nu$ into a conjugate, \unsure{are} inner, and that the center of $H_\nu$ is $\{1\}$.

% L5--9
In fact, $U_\nu$ can also be interpreted as a homogeneous space under $G_\nu^\circ$, and not only under $G_\nu$, more generally we have that $G_\nu^\circ$ is transitive on $U_\nu$, and \unsure{even} on $\mathcal P_{\nu'}(U_\nu)$ for every $\nu'\in\mathbb N^*$.

% L10--14, formule (30)
Returning to the model situation with $U_{g,\nu}$, we find
\[
\tag{30}
\begin{aligned}
U_{g,\nu}
&\simeq A_{g,\nu}^{\unsure{\deleted{\circ}|\bullet}}/
\underbrace{A_{g,(\nu,\nu+1)}}_{\substack{\text{\unsure{$\parallel$ definition}}\\ A_{g,\nu}\cap A_{g,\nu+1}}}\\[1.2em]
&\simeq A_{g,\nu}^{\circ}/
\underbrace{B_{g,(\nu,\nu+1)}}_{\substack{\text{\unsure{$\parallel$ (definition)}}\\
 A_{g,\deleted{\unsure{(\nu,}}\unsure{\{\nu+1\}}}\cap A_{g,\nu}^{\circ}
 =\{u\in A_{g,\nu}^{\circ}\mid u.a_{g,\nu}=\unsure{a_{g,\nu}}\}}}.
\end{aligned}
\]
% Exposant du premier numérateur : petite tache d'encre pleine (cercle noirci ou point), ni « ! » ni « * ».

% L15--18, bas de la page
We thus have \deleted{\unsure{some}} \deleted{\unsure{inclusions}} \inserted{by inverse} diagram of inclusions of subgroups of $A_{g,\nu}$ :
% « par inv. » est écrit en petit sous le mot barré « inclusions ».
\deleted{\diagram{\begin{tikzcd}[ampersand replacement=\&,column sep=1.7em,row sep=2.7em]
\unsure{A_{g,\illegible{1}}} \arrow[r,phantom,"\unsure{\supset}" description] \&
\unsure{A_{g,\nu}} \arrow[r,phantom,"\supset" description] \&
A_{g,\nu}^{!} \arrow[r,phantom,"\unsure{\supset}" description] \&
A_{g,\nu}^{\circ} \arrow[r,phantom,"\supset" description] \&
B_{g,\nu} \arrow[r,phantom,"\supset" description] \arrow[dl,hook,bend left=15,line width=0.8pt] \&
B_{g,\nu}^{\circ} \\
{} \& {} \& {} \& A_{g,(\nu,\nu+1)} \arrow[ull,hook,no head,bend left=10] \& {} \& {}
\end{tikzcd}}}
\margin{210}

\scanpage{212}
\mspage{106}
% L1--5, diagramme (31)
\diagram{\begin{tikzcd}[ampersand replacement=\&,column sep=2.1em,row sep=3.2em]
|[yshift=-1.4em]| {(31)} \&
A_{g,\nu} \arrow[rr,bend left=24,no head,"\Gamma_{g,\nu}"] \&
A_{g,\nu}^{!} \arrow[l,hook,"\mathrm{inv.}"',"\unsure{\mathfrak S_{\nu}}"] \&
A_{g,\nu}^{\circ} \arrow[l,hook,"\mathrm{inv.}"',"\Gamma_{g,\nu}^{!}"] \& {} \\
\left[A_{g,\nu+1}\,\hookleftarrow\right] \&
A_{g,(\nu,\nu+1)} \arrow[u,hook] \arrow[rr,bend right=45,no head,"\Gamma_{g,\nu}"'] \&
A_{g,\nu+1}^{!} \arrow[u,hook] \arrow[l,hook,"\substack{\mathrm{inv.}\\\unsure{\mathfrak S_{\nu}}}"] \&
B_{g,\nu} \arrow[u,hook] \arrow[l,hook,"\substack{\mathrm{inv.}\\\Gamma_{g,\nu}^{!}}"] \arrow[d,equal] \&
B_{g,\nu}^{\circ}=(A_{g,\nu+1}^{!})^{\circ}=A_{g,\nu+1}^{\circ} \arrow[l,hook,"\mathrm{inv.}"] \\[-1.2em]
{} \& {} \& {} \& \left[(A_{g,\nu}^{\circ})_{\unsure{a_{g,\nu}}}\right] \& {}
\end{tikzcd}}

% Marge gauche, à hauteur de L4--13
\margin{\uline{N.B.} $A_{g,(\nu,\nu+1)}$, \unsure{stabilizer} \illegible{2} in $(A_{g,\nu}=\operatorname{Aut}(U_{g,\nu}))$ of the \unsure{marked point} $a_{g,\nu}$ \unsure{of} \illegible{1}, \unsure{deserves} a simpler notation, for example $A_{g,\nu}^{\bullet}$ (the exponent $\bullet$ \unsure{here} to \unsure{emphasize} \illegible{3} \uline{\unsure{pointed|pointed}} in $U_{g,\nu}$).}

% L6--9
The two squares are cartesian, and the horizontal inclusions \inserted{(except the one in \unsure{brackets} in the $2^{\text{nd}}$ row)} are inclusions of invariant subgroups \inserted{invariant even} (in the largest group, $A_{g,\nu}$ resp. $A_{g,(\nu,\nu+1)}$).

% L10--15
[The equality $B_{g,\nu}^{\circ}=A_{g,\nu+1}^{!\,\circ}$ \deleted{\unsure{by}} (which refines the trivial inclusion $B_{g,\nu}^{\circ}\subset A_{g,\nu+1}^{!\,\circ}$ by the reverse inclusion $A_{g,\nu+1}^{!\,\circ}\subset B_{g,\nu}^{\circ}$, equivalent to $A_{g,\nu+1}^{!\,\circ}\subset B_{g,\nu}$) comes from the inclusion
\[
A_{g,\nu+1}^{!\,\circ}\inserted{\unsure{=}A_{g,\nu+1}^{\circ}}\subset A_{g,\nu}^{!\,\circ}=A_{g,\nu}^{\circ}
\quad\text{(whence : }A_{g,\nu+1}^{\circ}\subset A_{g,\nu+1}\cap A_{g,\nu}^{\circ}=B_{g,\nu}\text{).}
\]
% L'ajout interlinéaire est écrit au-dessus du premier \subset, relié à lui par un trait vertical que coupent deux petits traits.
]

% L16--19
The three vertical inclusions define three homogeneous spaces, and \deleted{$\unsure{A_g}$} inclusion homomorphisms between them which are \uline{bijective}
\[
(\unsure{U_{g,\nu}}\simeq)\quad
A_{g,\nu}/A_{g,(\nu,\nu+1)}
\xleftarrow{\sim} A_{g,\nu}^{!}/A_{g,\nu+1}^{!}
\xleftarrow{\sim} A_{g,\nu}^{\circ}/B_{g,\nu},
\]
% Le U de (U_{g,\nu}\simeq) est repassé à l'encre forte, non barré.
\inserted{\unsure{with} (transitivity !)}
and \deleted{\unsure{likewise by}} \deleted{\unsure{arrows} \illegible{2}} \inserted{the quotient groups \unsure{denoted} $\mathfrak S_\nu$ and $\Gamma_{g,\nu}^{!}$} defined by the horizontal inclusions of the first row, are isomorphic via the vertical inclusions to the corresponding quotients in the second row. Thus the extension of groups
\[
(12)\qquad 1\longrightarrow\Gamma_{g,\nu}^{!}\longrightarrow\Gamma_{g,\nu}\longrightarrow\unsure{\mathfrak S_\nu}\longrightarrow1
\]
% L25--26, fin de page
can \deleted{also} be deduced equally well from
\margin{211}

\scanpage{213}
% L1--5, formule (32)
the first row of (31), or from the first three terms of the second, in particular
\[
\tag{32}
\left\{\begin{aligned}
\unsure{\mathfrak S_\nu}&\simeq A_{g,(\nu,\nu+1)}/A_{g,\nu+1}^{!}\\
\Gamma_{g,\nu}^{\unsure{\bullet|\,}}&\simeq A_{g,(\nu,\nu+1)}/B_{g,\nu}\\
\Gamma_{g,\nu}^{!}&\simeq A_{g,\nu+1}^{!}/B_{g,\nu}.
\end{aligned}\right.
\]
% Au-dessus de \Gamma_{g,\nu} (deuxième ligne) : point d'encre plein, sans cercle ; peut-être une tache.

% L6--12
Thus the \inserted{\unsure{canonical}} torsor with group $A_{g,(\nu,\nu+1)}$ \inserted{$=(A_{g,\nu})_{\unsure{a_{g,\nu}}}$} over \inserted{the homogeneous space} $U_{g,\nu}$, and \inserted{the one} of the group $A_{g,\nu+1}^{!}$ ($\subset A_{g,(\nu,\nu+1)}$) which gives rise to it, is in either case deduced, by extension of the structure groups, from a torsor with group $B_{g,\nu}$ (whose fiber at $x\in U_{g,\nu}$ is the space of \deleted{$g$} $u\in A_{g,\nu}^{\circ}$ such that $u(a_{g,\nu})=x$).
\inserted{$\unsure{=(A_{g,\nu}^{\circ})_{a_{g,\nu}}}$}
% Entre « de groupe » et A_{g,(\nu,\nu+1)}, une expression (A_{g,\nu})_{a_{g,\nu}} fortement surchargée est écrite dans la ligne ; une autre, =(A_{g,\nu})_{a_{g,\nu}}, en haut à droite au-dessus de « sur ».

% L13--16
The associated \inserted{Galois} covering, with group $B_{g,\nu}/B_{g,\nu}^{\circ}$, is therefore also the homogeneous space $A_{g,\nu}^{\circ}/A_{g,\nu+1}^{\circ}$, which is obviously connected, and pointed over $a_{g,\nu}\in U_{g,\nu}$. Put
\deleted{\uline{\unsure{Theorem}}}
\deleted{\uline{\illegible{2}}}
\margin{212}
\[
\tag{33}
\widetilde U_{g,\nu}=A_{g,\nu}^{\circ}/A_{g,\nu+1}^{\circ}.
\]

% L17--21, théorème et ajouts
\uline{Theorem}. $\widetilde U_{g,\nu}$ is simply connected \inserted{(if $(g,\nu)\ne(1,0)$ and $(g,\nu)\ne(0,2)$)} (and even contractible if $(g,\nu)\ne(0,0)$) and thus identifies with the universal covering pointed at $\unsure{a_{g,\nu+1}}$ of $U_{g,\nu}$.

% Marge gauche, à hauteur de L17--22
\margin{\unsure{This relation|This assertion} \unsure{plausible} \illegible{3} in the anabelian cases $2g+\nu\geq3$, \unsure{we shall see} \illegible{3} in abelian case \illegible{4} \unsure{reasonable} \illegible{3}. We must exclude the cases $(1,0)$ and $(0,2)$.}

% L22--23, formule (34), au bord inférieur
\uline{Corollary}. We have a canonical isomorphism
\[
\tag{34}
B_{g,\nu}/B_{g,\nu}^{\circ}\simeq\pi_1(U_{g,\nu};a_{g,\nu})\overset{\text{definition}}{=}\Pi_{g,\nu}.
\]

\scanpage{214}
\mspage{107}
% L1--4
Proof of the Theorem, in terms of ``well-known'' things. The \deleted{\illegible{1}} homotopy exact sequence for the fibration of $A_{g,\nu}^{\circ}$ over
\[
A_{g,\nu+1}^{\circ}/A_{g,\nu}^{\circ}\simeq\widetilde U_{g,\nu}
\]
\deleted{\unsure{gives}}
% L5, suite (35)
\[
\tag{35}
\longrightarrow\pi_{i+1}(\widetilde U_{g,\nu})
\longrightarrow\pi_i(A_{g,\nu+1}^{\circ})
\longrightarrow\pi_i(A_{g,\nu\deleted{\unsure{+1}}}^{\circ})
\longrightarrow\pi_i(\widetilde U_{g,\nu})
\longrightarrow\pi_{i-1}(A_{g,\nu+1}^{\circ})
\]
and it shows that $\pi_1(\widetilde U_{g,\nu})$ is isomorphic to the cokernel of
\[
\pi_1(A_{g,\nu}^{\circ})\longrightarrow\pi_1(A_{g,\nu+1}^{\circ}),
\]
whose kernel is a quotient of $\pi_2(\widetilde U_{g,\nu})$\ldots

% L9--13, formule (36)
If $(g,\nu)\ne0$, we know that $\pi_i(\widetilde U_{g,\nu})=0$ for $i\geq2$ (for every covering of $U_{g,\nu}$) hence
\deleted{\illegible{1}}
\[
\tag{36}
\left\{\begin{gathered}
\inserted{\text{if }(g,\nu)\ne(0,0)}\\
\pi_i(A_{g,\nu+1}^{\circ})\longrightarrow\pi_i(A_{g,\nu\deleted{\unsure{+1}}}^{\circ})
\text{ is }\begin{cases}
\text{bijective if }i\geq2\\
\text{injective with invariant image if }i=1
\end{cases}\\
\text{[the cokernel is }\simeq\pi_1(\widetilde U_{g,\nu})\text{]}
\end{gathered}\right.
\]
and we want to prove that it is even bijective for $i=1$. The thing to remember \unsure{easy} is this :

% L17--25, ensemble barré en diagonale
\deleted{\uline{Well-known theorem}.
$\pi_1(A_{g,\nu}^{\circ})=1$ \unsure{iff} We are in the anabelian cases (i.e. $2g+\nu\geq3$), \deleted{\illegible{1}}
\inserted{($\pi_1(U_{g,\nu})$ nonabelian i.e.\illegible{1})}
\deleted{Thus} [\inserted{This is equivalent} \deleted{\illegible{2}}, taking account of (35), \unsure{to} \deleted{we have}
\[
\begin{aligned}
\pi_1(A_{g,0}^{\circ})&=1\quad\text{if }g\geq2\\
\pi_1(A_{1,1}^{\circ})&=1\\
\pi_1(A_{0,3}^{\circ})&=1
\end{aligned}
\qquad
\begin{gathered}
\pi_1(A_{1,0}^{\circ}),\pi_1(A_{0,0}^{\circ}),\pi_1(A_{0,1}^{\circ})\\
\text{and }\pi_1(A_{0,2}^{\circ})\text{ are }\ne1.
\end{gathered}
\]
]}
\margin{213}